\slajdid{04-s014}
\begin{frame}[label=04-s014]
\gusttekst
Primer prelaska\qquad $V_p=V_q$

Najbliže našem načinu razmišljanja\qquad $V_p\to V_{10}\to V_q$

Ne mora. Mi smo naviknuti od osnovnih škola na decimalni brojni sistem. Tablica sabiranja napamet. Tablica množenja napamet. I ti algoritmi računanja su nam postali urođeni, pa o njima i ne razmišljamo. Ali sada smo prinuđeni da bi „obučili“ naš digitalni sistem da se vratimo na ta računanja i ustanovimo algoritme.
\[236_7={?}_2\]
Računanje radimo u brojnom sistemu sa osnovom 7
\par\noindent\begin{tabularx}{\textwidth}{@{}lXXX@{}}
prvi korak&&$2{:}2=1\text{ ostatak }0$&$236_7\div2=1\ldots$\\[3mm]
drugi korak&$\textcolor{DEcrvena}{23}6_7\div2={?}$&$03{:}2=1\text{ ostatak }1$&$236_7\div2=11\ldots$\\[3mm]
treći korak&$\textcolor{DEcrvena}{236}_7\div2={?}$&$\textcolor{DEcrvena}{16{:}2=6\text{ ostatak }1}$&$236_7\div2=116$
\end{tabularx}\par
\end{frame}
